\chapter{مبرهنة فيثاغورس}\label{ch:g8:pythagoras}

أشهر مبرهنة في الرياضيات كلها تربط أضلاع \omterm{def:g6:shapes:triangles}{المثلث القائم
الزاوية} الثلاثة. وبها تصير أطوال قابلة للحساب لم تكن أيّ مسطرة
تستطيع قياسها مباشرةً: كالأقطار والمسافات والارتفاعات. وهي أيضًا
أول لقاء كبير لنا مع \omterm{ex:g1:subtraction:difference}{الفرق} بين المبرهنة
وعكسها.

\section{المبرهنة}

\begin{theorem}[فيثاغورس]\label{thm:g8:pythagoras:direct}
إذا كان المثلث $ABC$ \omterm{def:g6:shapes:triangles}{قائم الزاوية} في $A$، فإن
\[
BC^2 = AB^2 + AC^2 :
\]
أي إن مربع الوتر (وهو الضلع المقابل \omterm{def:g3:shapes:rightangle}{للزاوية القائمة})
يساوي \omterm{def:g1:addition:def}{مجموع} مربعَي الضلعين الآخرين.
\end{theorem}

\begin{proof}[فكرة البرهان]
أربع نسخ من المثلث، مرتّبة داخل مربع كبير طول ضلعه
$AB + AC$، تترك دون تغطية مربعًا مائلًا مبنيًّا على الوتر:
فمقارنة المساحات تعطي التساوي. والحساب المفصّل
مقترح في \cref{exo:g8:pythagoras:11}؛ والمبرهنة تنتج أيضًا
من \omterm{def:g2:mult:def}{الجداء} السلّمي، المدروس في مجلّد المرحلة الثانوية.
\end{proof}

\begin{omfigure}
\begin{tikzpicture}[scale=0.58]
  \coordinate (A) at (0,0);
  \coordinate (B) at (4,0);
  \coordinate (C) at (0,3);
  \draw[thick, fill=omDef!12] (A) -- (B) -- (C) -- cycle;
  \draw[thin] (0.35,0) -- (0.35,0.35) -- (0,0.35);
  \node[below left, font=\small] at (A) {$A$};
  \node[below right, font=\small] at (B) {$B$};
  \node[above left, font=\small] at (C) {$C$};
  % squares on the sides
  \draw[omDef, thick, fill=omDef!8] (0,0) rectangle (4,-4);
  \node[omDef, font=\small] at (2,-2) {$AB^2 = 16$};
  \draw[omMeth, thick, fill=omMeth!8] (0,0) rectangle (-3,3);
  \node[omMeth, font=\small] at (-1.5,1.5) {$AC^2 = 9$};
  \draw[omThm, thick, fill=omThm!8] (4,0) -- (7,4) -- (3,7) -- (0,3)
    -- cycle;
  \node[omThm, font=\small] at (3.5,3.5) {$BC^2 = 25$};
\end{tikzpicture}

{\small المبرهنة عبارةً عن مساحات: فالمربع المبني على
الوتر ($25$) له بالضبط \omterm{def:g6:measure:area}{مساحة} المربعين الآخرين
معًا ($16 + 9$). والأضلاع هنا $3$ و $4$ و $5$.}
\end{omfigure}

\begin{method}[حساب طول]\label{met:g8:pythagoras:compute}
في \omterm{def:g6:shapes:triangles}{مثلث قائم الزاوية} معلوم فيه ضلعان:
\begin{enumerate}
  \item سمِّ الوتر (وهو المقابل \omterm{def:g3:shapes:rightangle}{للزاوية القائمة} --- وأطول
        الأضلاع دائمًا)؛
  \item واُكتب تساوي \cref{thm:g8:pythagoras:direct} بالقيم
        المعلومة؛
  \item وحُلّ لإيجاد المربع الناقص: \emph{اِجمع} المربعين إذا كان
        الوتر مجهولًا، و\emph{اِطرح} إذا كان أحد ضلعَي القائمة مجهولًا؛
  \item وخذ الجذر التربيعي، بالضبط أو بالآلة الحاسبة، وتحقّق
        من أن الوتر خرج أطولها.
\end{enumerate}
\end{method}

\begin{example}[إيجاد الوتر]\label{ex:g8:pythagoras:hypotenuse}
\omterm{def:g6:shapes:triangles}{مثلث قائم الزاوية} ضلعا قائمته $6$ و $8$. والوتر $c$:
\[
c^2 = 6^2 + 8^2 = 36 + 64 = 100,
\qquad
c = \sqrt{100} = 10 .
\]
(والعدد $\sqrt{100}$، أي \emph{الجذر التربيعي} للعدد $100$، هو
العدد الموجب الذي مربعه $100$؛ وللجذور التربيعية فصل كامل
في \cref{ch:g9:sqrt}.)
\end{example}

\begin{example}[إيجاد ضلع قائمة]\label{ex:g8:pythagoras:leg}
سُلَّم طوله $5$ م مسنود إلى جدار، وقدمه على بُعد $1.4$ م من الجدار.
والارتفاع المبلوغ $h$: فالسُّلَّم هو الوتر، ومنه
\[
h^2 = 5^2 - 1.4^2 = 25 - 1.96 = 23.04,
\qquad
h = \sqrt{23.04} = 4.8 \text{ m}.
\]
اِطرح ولا تجمع: فالمجهول هنا \emph{ضلع قائمة}.
\end{example}

\section{المبرهنة العكسية}

\begin{theorem}[عكس فيثاغورس]\label{thm:g8:pythagoras:converse}
إذا كانت أضلاع مثلث $ABC$ تحقّق $BC^2 = AB^2 + AC^2$، فإن
\omterm{def:g6:shapes:triangles}{المثلث قائم الزاوية} في $A$.
\end{theorem}
\admitted

\begin{method}[اِختبار زاوية قائمة]\label{met:g8:pythagoras:test}
من أجل أضلاع مثلث معطاة:
\begin{enumerate}
  \item اِحسب على حدة مربع الضلع الأطول، \omterm{def:g1:addition:def}{ومجموع}
        مربعَي الضلعين الآخرين؛
  \item فإن تساوت النتيجتان، كان \omterm{def:g6:shapes:triangles}{المثلث قائم الزاوية}
        (بالعكس)، عند الرأس المقابل للضلع الأطول؛
  \item وإن اِختلفتا، فليس \omterm{def:g6:shapes:triangles}{قائم الزاوية} --- لأن مبرهنة
        فيثاغورس كانت ستفرض التساوي.
\end{enumerate}
\end{method}

\begin{example}\label{ex:g8:pythagoras:test}
الأضلاع $5$ و $12$ و $13$: مربع الضلع الأطول $13^2 = 169$؛ \omterm{def:g1:addition:def}{ومجموع}
المربعين الآخرين $5^2 + 12^2 = 25 + 144 = 169$. متساويان: إذن \omterm{def:g6:shapes:triangles}{قائم الزاوية}
(في مقابل الضلع $13$).

والأضلاع $6$ و $7$ و $9$: $9^2 = 81$ لكن $6^2 + 7^2 = 85 \neq 81$: فليس
مثلثًا \omterm{def:g6:shapes:triangles}{قائم الزاوية} (وهو «يكاد» يكون كذلك --- فالأعداد تقرّر، لا
العين).
\end{example}

\begin{remark}[المبرهنة وعكسها]
المبرهنة وعكسها يقولان أمرين مختلفين: فالمبرهنة
\emph{تستعمل} \omterm{def:g3:shapes:rightangle}{زاوية قائمة} لحساب أطوال؛ والعكس يستعمل
الأطوال \emph{ليبرهن} على \omterm{def:g3:shapes:rightangle}{زاوية قائمة}. وقد اِستعمل البنّاؤون
العكس منذ آلاف السنين: فحبل بعقد عند $3$ و $4$ و $5$ وحدات
مشدود يعطي \omterm{def:g3:shapes:rightangle}{زاوية قائمة} تامّة.
\end{remark}

\begin{example}[قطر المربع]\label{ex:g8:pythagoras:diagonal}
\omterm{def:g4:geometry:circle}{قطر} مربع طول ضلعه $1$ هو وتر \omterm{def:g6:shapes:triangles}{مثلث قائم
الزاوية} ضلعا قائمته $1$ و $1$:
\[
d^2 = 1^2 + 1^2 = 2 ,
\]
إذن $d = \sqrt2 \approx 1.414$ --- وهو عدد ليس كسرًا، كما
سيبرهن مجلّد المرحلة الثانوية. ومن أجل مربع طول ضلعه $c$، يكون
\omterm{def:g4:geometry:circle}{القطر} $c\sqrt2$.
\end{example}

\section{تمارين}

\begin{exercise}[$\star$]\label{exo:g8:pythagoras:1}
في كل \omterm{def:g6:shapes:triangles}{مثلث قائم الزاوية}، اِحسب الوتر:
ضلعا القائمة $9$ و $12$؛ \quad وضلعا القائمة $5$ و $12$؛ \quad وضلعا القائمة $8$ و $15$.
\end{exercise}

\begin{exercise}[$\star$]\label{exo:g8:pythagoras:2}
\omterm{def:g6:shapes:triangles}{مثلث قائم الزاوية} وتره $25$ وأحد ضلعَي قائمته $7$. اِحسب
ضلع القائمة الآخر.
\end{exercise}

\begin{exercise}[$\star$]\label{exo:g8:pythagoras:3}
اِحسب \omterm{def:g4:geometry:circle}{قطر} مستطيل ضلعاه $12$ سم و $9$ سم؛ \omterm{def:g4:geometry:circle}{وقطر}
مربع طول ضلعه $5$ سم (القيمة المضبوطة بجذر تربيعي،
ثم مقرّبة إلى المم).
\end{exercise}

\begin{exercise}[$\star$]\label{exo:g8:pythagoras:4}
أيّ المثلثات قائمة الزاوية؟ بَرِّر حسب
\cref{met:g8:pythagoras:test}:
\[
(20,\ 21,\ 29); \qquad
(7,\ 8,\ 11); \qquad
(1.5,\ 2,\ 2.5).
\]
\end{exercise}

\begin{exercise}[$\star$]\label{exo:g8:pythagoras:5}
بوابة مدعّمة بلوح قُطري. وعرض البوابة $3$ م
واِرتفاعها $1.6$ م: فكم يبلغ طول اللوح؟
\end{exercise}

\begin{exercise}[$\star$]\label{exo:g8:pythagoras:6}
\omterm{def:g6:shapes:triangles}{مثلث متساوي الساقين} ساقاه $10$ سم وقاعدته $12$ سم.
واِرتفاعه يشقّه إلى مثلثين قائمَي الزاوية: اِحسب ذلك الارتفاع،
ثم \omterm{def:g6:measure:area}{مساحة} المثلث.
\end{exercise}

\begin{exercise}[$\star\star$]\label{exo:g8:pythagoras:7}
شاشة تلفاز مستطيل $16{:}9$ قياسه $88.5$ سم في
$49.8$ سم. وتُباع الشاشات \omterm{def:g4:geometry:circle}{بالقطر}، بالبوصة
($1$ بوصة $= 2.54$ سم). اِحسب \omterm{def:g4:geometry:circle}{القطر} بالسنتيمتر، ثم بالبوصة
(مقرّبًا إلى أقرب بوصة).
\end{exercise}

\begin{exercise}[$\star\star$]\label{exo:g8:pythagoras:8}
يبحر قارب $24$ كم شرقًا، ثم $10$ كم شمالًا. فكم يبعد عن
نقطة اِنطلاقه (في خطّ \omterm{def:g6:lines:objects}{مستقيم})؟ اُرسم الوضعية أولًا.
\end{exercise}

\begin{exercise}[$\star\star$]\label{exo:g8:pythagoras:9}
سُلَّم طوله $2.5$ م يجب أن يبلغ عتبة نافذة على اِرتفاع $2.4$ م عن الأرض.
فكم يجب أن تبعد قدمه عن الجدار؟ وهل الجواب
متوافق مع نصيحة السلامة (القدم على نحو رُبع طول
السُّلَّم من الجدار)؟
\end{exercise}

\begin{exercise}[$\star\star$]\label{exo:g8:pythagoras:10}
على شبكة إحداثيات، ضع $A(1, 2)$ و $B(5, 5)$، واُرسم
\omterm{def:g6:shapes:triangles}{المثلث القائم الزاوية} الذي ضلعا قائمته موازيان للمحورين ووتره
$[AB]$. واِحسب $AB$. (وهذا الإنشاء، من أجل أيّ نقطتين، يصير
صيغة المسافة في الهندسة الإحداثية، في مجلّد المرحلة
الثانوية.)
\end{exercise}

\begin{exercise}[$\star\star\star$]\label{exo:g8:pythagoras:11}
أربع نسخ من \omterm{def:g6:shapes:triangles}{مثلث قائم الزاوية} ضلعا قائمته $a$ و $b$ ووتره $c$
موضوعة في أركان مربع طول ضلعه $a + b$، فتترك
مربعًا مائلًا طول ضلعه $c$ دون تغطية في الوسط.
\begin{enumerate}
  \item عبّر عن \omterm{def:g6:measure:area}{مساحة} المربع الكبير بطريقتين: مباشرةً، و
        على صورة (أربعة مثلثات) $+$ (المربع المائل).
  \item واُنشر $(a + b)^2$ (اُنظر \cref{thm:g8:equations:expand}) و
        اِستنتج $c^2 = a^2 + b^2$: فتكون قد برهنت على مبرهنة
        فيثاغورس.
\end{enumerate}
\end{exercise}

\section{مسألة: الثلاثيات الفيثاغورية}

\begin{problem}[{مسألة نهاية الأسبوع --- المثلثات القائمة الزاوية ذات الأعداد
الصحيحة، والمتطابقة $(m^2 - n^2)^2 + (2mn)^2 = (m^2 + n^2)^2$ التي
تولّدها كلها}]
\label{pb:g8:pythagoras:1}
\emph{الثلاثية الفيثاغورية}\index{ثلاثية فيثاغورية} ثلاثية
أعداد صحيحة $(a, b, c)$ تحقّق
\[
a^2 + b^2 = c^2 ,
\]
مثل ثلاثية البنّائين $(3, 4, 5)$. وقد فهرس لوح طيني بابلي
ثلاثيات ضخمة --- منها $(4601,\ 4800,\ 6649)$ --- قبل
فيثاغورس بأكثر من ألف سنة. وتصطاد هذه المسألة الثلاثيات،
ثم تبني الآلة التي تنتجها: متطابقة جبرية
مأخوذة مباشرةً من \cref{ch:g8:equations}، موضوعة في خدمة
الهندسة.

\medskip
\noindent\textbf{الجزء الأول --- صيد الثلاثيات.}
\begin{enumerate}
  \item تحقّق من أن $(5, 12, 13)$ و $(8, 15, 17)$ و
        $(20, 21, 29)$ ثلاثيات فيثاغورية. فماذا تقول
        \cref{thm:g8:pythagoras:converse} عن المثلثات ذات
        أطوال الأضلاع هذه؟
  \item اِشرح لماذا يعطي حبل مشدود ذو اِثنتي عشرة \omterm{def:g6:lines:objects}{قطعة} متساوية، موضوع
        مثلثًا أضلاعه $3$ و $4$ و $5$ قطع،
        البنّائين \omterm{def:g3:shapes:rightangle}{زاوية قائمة} تامّة.
  \item بيّن أنه إذا كانت $(a, b, c)$ \omterm{pb:g8:pythagoras:1}{ثلاثية فيثاغورية}، فكذلك
        $(ka, kb, kc)$ من أجل كل عدد صحيح $k \geq 1$. واِستنتج
        ثلاث ثلاثيات جديدة من $(3, 4, 5)$.
  \item \omterm{def:g5:division:multiple}{مضاعفات} $(3, 4, 5)$ هي الثلاثيات
        $(3k, 4k, 5k)$. بيّن أن $(5, 12, 13)$ \emph{ليست} واحدة
        منها.
  \item اِشرح لماذا يبرهن السؤالان 3 و4 معًا على أن تكبير
        $(3, 4, 5)$ لا يمكن أن ينتج كل \omterm{pb:g8:pythagoras:1}{ثلاثية فيثاغورية}:
        فلا بدّ من آلة أخرى.
\end{enumerate}

\medskip
\noindent\textbf{الجزء الثاني --- الآلة.}
خذ عددين صحيحين $m > n \geq 1$ وكوّن الأعداد الثلاثة
\[
a = m^2 - n^2,
\qquad
b = 2mn,
\qquad
c = m^2 + n^2 .
\]
\begin{enumerate}[resume]
  \item باستعمال المتطابقات الشهيرة في
        \cref{pb:g8:equations:1} (اُكتب $A = m^2$ و $B = n^2$)،
        اُنشر $a^2$ و $b^2$، وبَرهِن على أن
        \[
        \left(m^2 - n^2\right)^2 + (2mn)^2
        = \left(m^2 + n^2\right)^2 :
        \]
        أي إن الآلة تُخرج دائمًا \omterm{pb:g8:pythagoras:1}{ثلاثية فيثاغورية}.
  \item شغّل الآلة على $(m, n) = (2, 1)$ و $(3, 1)$ و $(3, 2)$
        و $(4, 1)$ و $(4, 3)$، واِعرض الثلاثيات الخمس في
        جدول.
  \item إحدى الثلاثيات في جدولك نسخة مكبّرة من أصغر
        منها. فأيّها، ومن أيّها؟
  \item جِد $(m, n)$ التي تُخرج بها الآلة الثلاثية
        $(20, 21, 29)$ في السؤال 1 (بضلعَي قائمة بأيّ ترتيب).
  \item جِد $m$ و $n$ يحقّقان $m^2 + n^2 = 100$، واِستنتج
        \omterm{pb:g8:pythagoras:1}{ثلاثية فيثاغورية} وترها $100$.
\end{enumerate}

\medskip
\noindent\textbf{الجزء الثالث --- العائلات، وخاتمة زوجية.}
\begin{enumerate}[resume]
  \item شغّل الآلة مع $m = n + 1$. بيّن أن ضلع القائمة الفردي
        هو $2n + 1$، وضلع القائمة الزوجي $2n^2 + 2n$، والوتر
        $2n^2 + 2n + 1$ --- إذن يفترق الوتر وضلع القائمة الزوجي بمقدار
        $1$ بالضبط. واُكتب الثلاثيات من أجل $n = 1, 2, 3, 4$.
  \item اِستنتج أن \emph{كل} \omterm{def:g3:numbers:evenodd}{عدد فردي} $3, 5, 7, 9, \dots$
        هو ضلع \omterm{def:g6:shapes:triangles}{لمثلث قائم الزاوية} ذي أعداد صحيحة، وأعطِ
        ثلاثية تحوي $11$.
  \item \omterm{def:g3:numbers:evenodd}{والأعداد الزوجية} تصلح أيضًا: شغّل الآلة مع $n = 1$، و
        أنتج ثلاثية تحوي $14$. ثم اِحصل على الثلاثية
        نفسها مرة ثانية، بتكبير واحدة من جدول
        السؤال 7.
  \item بيّن أن مربع \omterm{def:g3:numbers:evenodd}{العدد الزوجي} زوجي وأن
        مربع \omterm{def:g3:numbers:evenodd}{العدد الفردي} فردي. (اُكتب العدد على صورة
        $2k$ أو $2k + 1$ واِستعمل متطابقة.)
  \item اُختم بخاتمة الزوجية: لا توجد \omterm{pb:g8:pythagoras:1}{ثلاثية فيثاغورية}
        أعدادها الثلاثة كلها فردية --- ففي كل \omterm{def:g6:shapes:triangles}{مثلث قائم الزاوية}
        ذي أعداد صحيحة، ضلع واحد على الأقلّ زوجي.
\end{enumerate}
\end{problem}
